\documentclass[10pt,a4paper]{amsart}
\usepackage{amsmath,amssymb,amsthm,graphicx}
\usepackage{amsmath}
\usepackage{amssymb}
\usepackage{hyperref}
\newtheorem{Theorem}{Theorem}[section]
\newtheorem{Corollary}[Theorem]{Corollary}
\newtheorem{Lemma}[Theorem]{Lemma}
\newtheorem{Proposition}[Theorem]{Proposition}
{\theoremstyle{definition}
\newtheorem{Note}[Theorem]{Note}
\newtheorem{Definition}[Theorem]{Definition}
\newtheorem{Problem}[Theorem]{Problem}
}
\usepackage{cleveref}
\crefname{originaldefinition}{original definition}{original definitions}
\Crefname{originaldefinition}{Original definition}{Original definitions}
\makeatletter
\newlabel{def:abc}{{1.3}{source}{}{sourcepointer-1}{}}
\newlabel{def:abc@cref}{{[originaldefinition][0][]1.3}{[1][0][]source}{}{}{}}
\makeatother
\title{Six explicit words in the Askey-Wilson generators define one and the same element}
\author{Paul Terwilliger}
\date{doi:10.3842/SIGMA.2011.069; CC BY4.0}
\begin{document}
\maketitle
\noindent\textit{Editorial note.} All delivered mathematical text is the author's own, copied from the pinned publisher source sources/sigma-2011-069-askey-wilson/pinned-original-source.tex; no statement or proof step was rewritten, reordered or removed.\par\medskip
\noindent\textit{Editorial note.} Preamble provenance: the retained preamble is the source's own layout adapter plus the author's macro block at lines 10-18 of sources/sigma-2011-069-askey-wilson/pinned-original-source.tex (everything between the publisher class metadata and \textbackslash{}begin\{document\}, comments excluded). Nothing was added, no macro was invented, and the publisher document class is deliberately not reproduced; the excerpt serves only as preamble, never as a span coordinate.\par\medskip
\noindent\textit{Editorial note.} External references are navigation pointers only; every deliberately kept step of the author's argument stays inside the counted spans. The displayed numbers of those pointers are the author's own numbering under the author's \textbackslash{}newtheorem declarations (the source has no subsections, so numbering is section.index), computed mechanically by the lane's scanner.\par\medskip
\noindent\textit{Editorial note.} The only reference to another labelled item is Definition def:abc, recorded as a navigation stub of kind definition with the author's own number; every step that uses it stays inside the counted span.\par\medskip
\noindent\textit{Editorial note.} The counted spans contain no bibliography citation, so no bibliography excerpt is attached to this card.\par\medskip
\noindent\textit{Editorial note.} Licence binding: the card's source\_version now carries the value recorded in this source's own licence evidence (sources/sigma-2011-069-askey-wilson/license-observation.json), namely 'SIGMA7(2011),069; publisher final form2011-07-09'. The adopt step of the publication line compares the two and refused the earlier revision; the card is aligned to the evidence rather than the evidence to the card.\par\medskip
\noindent\textit{Allowed standard prerequisites (editorial).}\begin{itemize}
\item The defining relations of the universal Askey-Wilson algebra, arranged as the relations involving alpha, beta, gamma. Same source, Definition labelled def:abc; invoked inside the counted proof at lines 1511-1512 and 1526-1527. Supplies the two relations whose vanishing makes the displayed differences zero.
\item The automorphism rho of the universal Askey-Wilson algebra. Same source, set up before the statement; used inside the counted proof at lines 1495-1501. Cyclically permutes the six displayed elements, which reduces the verification to two computations.
\end{itemize}\medskip
\ifdefined\hypertarget\hypertarget{sourcepointer-1}{}\fi
\noindent\textit{Source reference.} Original definition 1.3: sources/sigma-2011-069-askey-wilson/pinned-original-source.tex, source lines 1512--1512. Navigation pointer for \textbackslash{}ref\{def:abc\} used inside the counted span (line 1512); the argument that consumes it is retained in the span.\par\medskip
\begin{Lemma}
\label{lem:six}
The following elements of $\Delta $ coincide:
\begin{gather*}
qABC + q^2A^2+q^{-2}B^2+q^2C^2-q A\alpha -q^{-1} B\beta -q C\gamma,
\\
qBCA +
q^2A^2 +
q^2B^2+q^{-2}C^2
-q A\alpha
-qB\beta -q^{-1}C\gamma,
\\
qCAB +
q^{-2}A^2+q^2B^2
+q^2 C^2
-q^{-1} A\alpha -q B \beta
-q C\gamma,
\\
q^{-1}CBA + q^{-2}A^2+q^{2}B^2+q^{-2}C^2-q^{-1}A\alpha -q B\beta -
q^{-1} C\gamma,
\\
q^{-1}ACB +
q^{-2}A^2 +
q^{-2}B^2+q^{2}C^2
-q^{-1} A\alpha
-q^{-1}B\beta -q C\gamma,
\\
q^{-1}BAC +
q^{2}A^2+q^{-2}B^2
+q^{-2}C^2
-q A\alpha
-q^{-1} B \beta
-q^{-1} C\gamma.
\end{gather*}
We denote this common element by $\Omega$.
\end{Lemma}

\begin{proof}
Denote the displayed sequence of elements by
$\Omega^+_B$,
$\Omega^+_C$,
$\Omega^+_A$,
$\Omega^-_B$,
$\Omega^-_C$,
$\Omega^-_A$.
The automorphism
$\rho$ cyclically permutes
$\Omega^+_A$,
$\Omega^+_B$,
$\Omega^+_C$
and cyclically permutes
$\Omega^-_A$, $\Omega^-_B$, $\Omega^-_C$.
The element
$\Omega^+_B-\Omega^-_C$
is
equal to $(q-q^{-1})A$ times
\begin{gather}
\label{eq:need0}
\big(q+q^{-1}\big)A + \frac{qBC-q^{-1}CB}{q-q^{-1}} -\alpha.
\end{gather}
The element
(\ref{eq:need0}) is zero by Def\/inition
\ref{def:abc} so
$\Omega^+_B=\Omega^-_C$.
In this equation we apply $\rho$ twice to get
$\Omega^+_C=\Omega^-_A$
and
$\Omega^+_A=\Omega^-_B$.
The element
$\Omega^+_B-\Omega^-_A$ is equal to
\begin{gather}
\label{eq:need02}
\big(q+q^{-1}\big)C + \frac{qAB-q^{-1}BA}{q-q^{-1}} -\gamma
\end{gather}
times $(q-q^{-1})C$.
The element
(\ref{eq:need02}) is zero by Def\/inition
\ref{def:abc} so
$\Omega^+_B=\Omega^-_A$.
Applying $\rho$ twice we get
$\Omega^+_C=\Omega^-_B$
and
$\Omega^+_A=\Omega^-_C$.
By these comments
$\Omega^+_B$,
$\Omega^+_C$,
$\Omega^+_A$,
$\Omega^-_B$,
$\Omega^-_C$,
$\Omega^-_A$ coincide.
\end{proof}

\begin{thebibliography}{99}

\footnotesize\itemsep=0pt
\bibitem{alpern}
Alperin R.C.,
Notes: ${\rm  {PSL}}_2(\mathbb Z)= \mathbb Z_2 \star \mathbb Z_3$,
\href{http://dx.doi.org/10.2307/2324963}{{\it Amer. Math. Monthly}} {\bf 100} (1993), 385--386.
\end{thebibliography}
\end{document}
